\newpage
\fancyhead[L]{\songti\zihao{-5}\cquHeader}
\fancyhead[R]{\songti\zihao{-5}附录F：通过 Comparator 进行声明验证}
\addcontentsline{toc}{section}{附录F：通过 Comparator 进行声明验证}
\section*{附录F：通过 Comparator 进行声明验证}
\zihao{5}
\begingroup
\setlength{\parindent}{0pt}
\setlength{\parskip}{0.35\baselineskip}

如第 3 章所述，本文从两个层面验证形式化结果。首先，整个 Lean 项目通过 \texttt{lake build}，确认项目中的声明均被 Lean 接受，且维护的源文件中不含未完成的证明义务。其次，本文使用 Comparator 工具~\cite{Comparator2026}，将简短的可信规范与完整项目中的定理声明进行比较，并把导出的环境重新交给 Lean 内核检查。

验证时临时放置在 \path{Mordell/Challenge.lean} 的规范模块如下所示。它只包含本文讨论的 Mordell 方程最终结论，而不包含证明细节。人工审阅者无需阅读 \texttt{Descent.lean}、\texttt{ClassGroupApprox.lean} 或具体类数文件中的证明工程，即可核对这些声明本身的语义是否正确。随后，Comparator 自动检查完整项目中同名声明是否证明了完全相同的命题，并且只依赖允许的 Lean 公理。

\begin{leanlisting}
import Mathlib.Tactic

theorem mordell_minus1_solution_iff (x y : Int) :
    y ^ 2 = x ^ 3 - 1 <-> x = 1 /\ y = 0 := by
  sorry

theorem mordell_minus2_solutions_iff (x y : Int) :
    y ^ 2 = x ^ 3 - 2 <->
      (x = 3 /\ y = 5) \/ (x = 3 /\ y = -5) := by
  sorry

theorem mordell_minus5 (x y : Int) :
    Not (y ^ 2 = x ^ 3 - 5) := by
  sorry

theorem mordell_minus6 (x y : Int) :
    Not (y ^ 2 = x ^ 3 - 6) := by
  sorry

theorem mordell_minus13_solutions_iff (x y : Int) :
    y ^ 2 = x ^ 3 - 13 <->
      (x = 17 /\ y = 70) \/ (x = 17 /\ y = -70) := by
  sorry
\end{leanlisting}

对应的 Comparator 配置如下：

\begin{leanlisting}
{
  "challenge_module": "Mordell.Challenge",
  "solution_module": "Mordell",
  "theorem_names": [
    "mordell_minus1_solution_iff",
    "mordell_minus2_solutions_iff",
    "mordell_minus5",
    "mordell_minus6",
    "mordell_minus13_solutions_iff"
  ],
  "permitted_axioms": [
    "propext",
    "Quot.sound",
    "Classical.choice"
  ],
  "enable_nanoda": false
}
\end{leanlisting}

运行 Comparator 时需要进入 Lake 环境，从而保证 \texttt{lean4export}、项目库路径以及锁定的 Lean 工具链都来自同一份构建配置：

\begin{leanlisting}
lake build Mordell
lake build comparator
lake build lean4export
lake env comparator comparator-mordell.json
\end{leanlisting}

在本文使用的开发环境中，项目固定在 Lean \texttt{v4.30.0-rc2}。完整项目可以成功构建。对 Mordell 最终解集声明运行 Comparator 后，得到如下结果：

\begin{leanlisting}
Running Lean default kernel on solution.
Lean default kernel accepts the solution
Your solution is okay!
\end{leanlisting}

该输出表明，导出的声明被 Lean 内核接受，并且被检查定理的证明体没有使用超过 \texttt{propext}、\texttt{Quot.sound} 和 \texttt{Classical.choice} 的额外公理。
\endgroup
